% AA-21 | Algebra abstrakcyjna | Iteracja XI
\chapter{Zanurzenia ciał, automorfizmy i rozszerzenia Galois}
\label{aa:ch:embeddings-normal-separable}

\section{Dlaczego pierwiastki można przestawiać?}
W Rozdziale~\ref{aa:ch:field-extensions} skonstruowaliśmy ciała
rozkładu wielomianów. Rozpatrzmy $L=\mathbb Q(\sqrt2)$.
Każdy element ma jednoznacznie postać $a+b\sqrt2$, a przyporządkowanie
$a+b\sqrt2\mapsto a-b\sqrt2$ zachowuje dodawanie, mnożenie oraz
wszystkie liczby wymierne. Nie jest to przypadek: każdy homomorfizm
ciał zachowuje równania wielomianowe, zatem musi wysyłać pierwiastek
wielomianu na inny pierwiastek tego samego wielomianu. Nie wszystkie
takie zamiany są jednak możliwe wewnątrz rozważanego ciała.

W całym rozdziale wszystkie ciała są przemienne, a ich homomorfizmy
zachowują jedynkę. Dla stosowania zanurzeń wygodnie jest pracować
w ciele, w którym każdy wielomian ma pierwiastek. Następująca
konstrukcja uzasadnia istnienie takiego rozszerzenia bez zakładania
go jako nieudowodnionego postulatu.

\section{Istnienie domknięcia algebraicznego}
\begin{definicja}[Ciało algebraicznie domknięte i domknięcie algebraiczne]
\label{aa:def:algebraic-closure}
Ciało $\Omega$ jest \emph{algebraicznie domknięte}, jeśli każdy
wielomian dodatniego stopnia z $\Omega[X]$ ma pierwiastek w $\Omega$.
\emph{Domknięciem algebraicznym} ciała $K$ nazywamy rozszerzenie
algebraiczne $\Omega/K$, w którym $\Omega$ jest algebraicznie domknięte.
\end{definicja}

\begin{twierdzenie}[Istnienie domknięcia algebraicznego]
\label{aa:tw:algebraic-closure-exists}
Każde ciało $K$ ma domknięcie algebraiczne.
\end{twierdzenie}
\begin{proof}
\emph{Krok 1: jednoczesne dołączenie pierwiastków wielomianów nad $K$.}
Niech $\mathcal P$ oznacza zbiór wszystkich unormowanych
wielomianów dodatniego stopnia z $K[X]$. Dla każdego
$f\in\mathcal P$ wprowadźmy osobną zmienną $Y_f$.
W pierścieniu wielomianów
$A=K[Y_f:f\in\mathcal P]$ rozważmy ideał
$I=(f(Y_f):f\in\mathcal P)$.
Wykazujemy, że $I$ jest właściwy.
Gdyby $1\in I$, istniałaby skończona relacja
$1=\sum_{j=1}^r h_jf_j(Y_{f_j})$ z $h_j\in A$.
Na mocy konstrukcji przez kolejne dołączanie
pierwiastków (Twierdzenie~\ref{aa:tw:adjoin-root-abstract})
możemy kolejno dołączać pierwiastki $f_1,\ldots,f_r$.
W każdym kroku otrzymujemy rozszerzenie skończone,
a twierdzenie o wieży zapewnia, że wynikowe $E/K$
również jest skończone. Wybierzmy w $E$ pierwiastek
$\alpha_j$ każdego $f_j$. Możemy zdefiniować homomorfizm
$\operatorname{ev}:A\to E$, podstawiając
$Y_{f_j}\mapsto\alpha_j$ i pozostałe zmienne
na przykład na $0$; jeśli dany wielomian występuje
wiele razy, wybieramy dla niego jedną wartość.
Po zastosowaniu $\operatorname{ev}$ otrzymalibyśmy
$1=\sum_j\operatorname{ev}(h_j)f_j(\alpha_j)=0$.
Jest to sprzeczność, więc $I\ne A$.

Rozważmy rodzinę wszystkich ideałów właściwych pierścienia $A$
zawierających $I$, uporządkowaną przez inkluzję. Jest ona
niepusta, a suma mnogościowa każdego niepustego łańcucha
z tej rodziny jest ideałem właściwym: gdyby zawierała $1$, element $1$
należałby już do jednego z ideałów łańcucha. Zatem
lemat Kuratowskiego--Zorna daje element maksymalny tej
rodziny (pusty łańcuch jest ograniczony przez $I$). Taki ideał $\mathfrak m$ jest maksymalnym ideałem
$A$: każdy ideał większy od $\mathfrak m$ zawiera $I$,
więc nie może być właściwy. W konsekwencji iloraz
$E=A/\mathfrak m$ jest ciałem.
Ponieważ $I$ jest właściwy, homomorfizm
$K\to E$ jest różnowartościowy: jego jądro jest
ideałem ciała $K$ i nie może zawierać $1$.
Dla każdego $f\in\mathcal P$ obraz $y_f$ zmiennej
$Y_f$ w $E$ spełnia $f(y_f)=0$.
Ciało $E=A/\mathfrak m$ jest generowane jako algebra nad $K$
przez elementy $y_f$, ponieważ każdy jego element
jest obrazem wielomianu z $A$. Każdy $y_f$ jest
algebraiczny nad $K$, a obraz dowolnego wielomianu
z $A$ zależy tylko od skończenie wielu zmiennych.
Zatem każdy $x\in E$ należy do
$K(y_{f_1},\ldots,y_{f_s})$, skończonego rozszerzenia
$K$ na mocy wielokrotnego zastosowania twierdzenia
o wieży. Każdy taki $x$ jest więc algebraiczny nad $K$.
Otrzymaliśmy rozszerzenie algebraiczne $E/K$, w którym
każdy niezerowy wielomian dodatniego stopnia z $K[X]$
ma pierwiastek.

\emph{Krok 2: iteracja konstrukcji.}
Zaczynamy od $K_0=K$. Dla każdego $n\ge0$ stosujemy
Krok 1 do $K_n$ i otrzymujemy rozszerzenie algebraiczne
$K_{n+1}/K_n$ takie, że każdy wielomian dodatniego
stopnia z $K_n[X]$ ma pierwiastek w $K_{n+1}$.
Utożsamiając kolejne ciała z podciałami większych,
określamy $\Omega=\bigcup_{n\ge0}K_n$.
Jest ono ciałem: każde skończenie wiele jego elementów
należy do wspólnego $K_n$, w którym można wykonywać
działania i brać odwrotności. Rozszerzenie $\Omega/K$
jest algebraiczne, gdyż każdy $x\in\Omega$ należy
do pewnego $K_n$. Uzasadnijmy używaną tu
przechodniość algebraiczności. Jeśli $B/A$ jest
algebraiczne, a $x$ jest algebraiczny nad $B$,
to współczynniki wielomianu, który zeruje $x$,
należą do pewnego skończonego rozszerzenia
$A(b_0,\ldots,b_d)/A$. Zatem $x$ należy do skończonego
rozszerzenia $A(b_0,\ldots,b_d,x)/A$, na mocy
twierdzenia o wieży, i jest algebraiczny nad $A$.
Stosując to skończenie wiele razy do wieży
$K_0\subseteq K_1\subseteq\cdots\subseteq K_n$,
widzimy, że $x$ jest algebraiczny nad $K$.

Niech teraz $f\in\Omega[X]$ ma dodatni stopień.
Wszystkie jego współczynniki należą do pewnego $K_n$,
więc $f$ ma pierwiastek w $K_{n+1}\subseteq\Omega$.
Zatem $\Omega$ jest algebraicznie domknięte i stanowi
domknięcie algebraiczne $K$.
W powyższej konstrukcji istnienie ideału maksymalnego wynika
z lematu Kuratowskiego--Zorna; bez tego kroku nie można zakładać,
że iloraz przez dowolny ideał będzie ciałem.
\end{proof}

W dalszej części przez $\Omega$ oznaczamy odpowiednie
domknięcie algebraiczne zawierające badane skończone
rozszerzenie. Konstrukcja przez wybór kolejnych pierwiastków
pozwala też umieścić każde ustalone rozszerzenie algebraiczne
w takim ciele. Wszystkie użyte niżej przedłużenia zanurzeń
dotyczą rozszerzeń skończonych.

\begin{definicja}[Zanurzenie nad ciałem bazowym i automorfizm]
\label{aa:def:field-k-embedding}
Niech $L/K$ oraz $M/K$ będą rozszerzeniami ciał.
\emph{$K$-zanurzeniem} $L$ w $M$ nazywamy homomorfizm
ciał $\sigma:L\to M$ taki, że $\sigma(a)=a$ dla każdego $a\in K$.
\emph{$K$-automorfizmem} $L$ nazywamy bijektywne $K$-zanurzenie
$L\to L$. Zbiór tych automorfizmów oznaczamy
$\operatorname{Aut}_K(L)$. Jest on grupą względem składania:
składanie i odwrotność bijektywnych homomorfizmów ciał
ponownie zachowują działania, jedynkę i każdy element $K$.
Jeśli $\varphi:K\to K'$ jest dowolnym zanurzeniem, to
$\varphi$ przenosi wielomian $f=\sum a_iX^i\in K[X]$
na $f^\varphi=\sum\varphi(a_i)X^i\in K'[X]$.
\end{definicja}

\begin{stwierdzenie}[Automorfizm zachowuje wielomian minimalny]
\label{aa:st:automorphism-minimal-polynomial}
Jeśli $L/K$ jest rozszerzeniem algebraicznym,
$\sigma\in\operatorname{Aut}_K(L)$ i $\alpha\in L$, to
$m_{\sigma(\alpha),K}=m_{\alpha,K}$.
W szczególności $m_{\alpha,K}(\sigma(\alpha))=0$,
a więc $\sigma(\alpha)$ jest jednym z pierwiastków
$m_{\alpha,K}$ leżących w $L$.
\end{stwierdzenie}
\begin{proof}
Dla $f(X)=\sum_{j=0}^d a_jX^j\in K[X]$ mamy
$\sigma(f(\alpha))=\sum a_j\sigma(\alpha)^j=f(\sigma(\alpha))$,
ponieważ $\sigma(a_j)=a_j$. Stosujemy tę równość do
$f=m_{\alpha,K}$, dla którego $f(\alpha)=0$.
Ponadto $m_{\alpha,K}$ jest unormowany i nierozkładalny,
więc jest także wielomianem minimalnym $\sigma(\alpha)$:
oba wielomiany są unormowane, a minimalny dzieli
$m_{\alpha,K}$.
\end{proof}

\begin{przyklad}[Dwa zanurzenia rozszerzenia kwadratowego]
\label{aa:prz:quadratic-conjugation}
Niech $L=\mathbb Q(\sqrt2)$. Są dokładnie dwa
$\mathbb Q$-automorfizmy $L$: identyczność i
$\tau(a+b\sqrt2)=a-b\sqrt2$. Istotnie, obraz $\sqrt2$
musi należeć do $\{\sqrt2,-\sqrt2\}$, a ten obraz jednoznacznie
wyznacza homomorfizm. Wzór na $\tau$ zachowuje mnożenie, ponieważ
$(a+b\sqrt2)(c+d\sqrt2)=(ac+2bd)+(ad+bc)\sqrt2$.
Wynika stąd $\operatorname{Aut}_{\mathbb Q}(L)\cong C_2$.
\end{przyklad}

\section{Przedłużanie zanurzeń}
\begin{lemat}[Przedłużenie zanurzenia na proste rozszerzenie]
\label{aa:lem:extend-embedding-simple}
Niech $\varphi:K\to K'$ będzie izomorfizmem ciał,
niech $\alpha$ będzie algebraiczny nad $K$, a
$m=m_{\alpha,K}\in K[X]$ jego wielomianem minimalnym.
Jeżeli $M/K'$ jest rozszerzeniem oraz $\beta\in M$ spełnia
$m^\varphi(\beta)=0$, to istnieje dokładnie jedno zanurzenie
$\widetilde\varphi:K(\alpha)\to M$ przedłużające
$\varphi$ i spełniające $\widetilde\varphi(\alpha)=\beta$.
\end{lemat}
\begin{proof}
Z Twierdzenia~\ref{aa:tw:simple-algebraic-extension} mamy
$K(\alpha)\cong K[X]/(m)$. Na przedstawicielach definiujemy
$\widetilde\varphi([f])=f^\varphi(\beta)$.
Jeśli $[f]=[g]$, to $f-g=mh$ dla pewnego $h\in K[X]$;
stąd $f^\varphi(\beta)-g^\varphi(\beta)=
m^\varphi(\beta)h^\varphi(\beta)=0$.
Zatem wzór jest poprawnie określony. Zachowanie działań
i jedynki wynika z tożsamości
$(f+g)^\varphi=f^\varphi+g^\varphi$ oraz
$(fg)^\varphi=f^\varphi g^\varphi$ i z własności
podstawiania wielomianów. Jądro homomorfizmu z ciała
jest ideałem, a ponieważ $1$ przechodzi na $1$, nie
jest całym ciałem. Zatem jądro jest zerowe i odwzorowanie
jest różnowartościowe; ponadto przyjmuje żądane wartości
na $K$ i $\alpha$.
Każdy element $K(\alpha)$ jest postaci $f(\alpha)$ dla
pewnego $f\in K[X]$, więc wartości na $K$ i $\alpha$
wyznaczają całe odwzorowanie jednoznacznie.
\end{proof}

\begin{wniosek}[Przedłużenie zanurzenia na rozszerzenie skończone]
\label{aa:wn:extend-embedding-finite}
Niech $L/K$ będzie rozszerzeniem skończonym,
$\varphi:K\hookrightarrow\Omega$ zanurzeniem w ciało
algebraicznie domknięte $\Omega$. Wtedy istnieje zanurzenie
$\widetilde\varphi:L\hookrightarrow\Omega$ przedłużające
$\varphi$. To samo zachodzi dla zadanego zanurzenia dowolnego
ciała pośredniego $E$, $K\subseteq E\subseteq L$.
\end{wniosek}
\begin{proof}
Ponieważ $[L:K]<\infty$, możemy wybrać skończony zbiór
generatorów $L=K(\alpha_1,\ldots,\alpha_r)$, na przykład
dowolną bazę przestrzeni wektorowej $L$ nad $K$.
Budujemy kolejne ciała $K_i=K(\alpha_1,\ldots,\alpha_i)$.
Jeżeli zanurzenie $\varphi_{i-1}:K_{i-1}\hookrightarrow\Omega$
już skonstruowano, to wielomian minimalny $\alpha_i$
nad $K_{i-1}$ po przeniesieniu współczynników przez
$\varphi_{i-1}$ ma pierwiastek w $\Omega$. Lemat
\ref{aa:lem:extend-embedding-simple}, stosowany z
$K'=\varphi_{i-1}(K_{i-1})$, daje jego przedłużenie
na $K_i$. Po skończonej liczbie kroków dostajemy
zanurzenie całego $L$. Dla ciała pośredniego $E$
powtarzamy argument, biorąc bazę $L$ nad $E$;
stopień $[L:E]$ jest skończony na mocy
Twierdzenia~\ref{aa:tw:field-tower}.
\end{proof}

\begin{twierdzenie}[Liczba zanurzeń i rozdzielczość]
\label{aa:tw:embeddings-count}
Niech $L/K$ będzie rozszerzeniem skończonym,
a $\Omega$ ciałem algebraicznie domkniętym zawierającym
$L$. Liczba różnych $K$-zanurzeń $L\hookrightarrow\Omega$
nie przekracza $[L:K]$. Jest równa $[L:K]$ wtedy
i tylko wtedy, gdy każdy element $L$ jest rozdzielczy
nad $K$, czyli gdy $L/K$ jest rozszerzeniem rozdzielczym.
\end{twierdzenie}
\begin{proof}
Zapiszmy $L=K(\alpha_1,\ldots,\alpha_r)$ oraz
$K_i=K(\alpha_1,\ldots,\alpha_i)$, $K_0=K$.
Oznaczmy $d_i=[K_i:K_{i-1}]$.
Gdy wybrano zanurzenie $\varphi_{i-1}:K_{i-1}\to\Omega$,
jego przedłużenie na $K_i$ jest jednoznacznie wyznaczone
przez obraz $\alpha_i$. Wskazany obraz musi być
pierwiastkiem wielomianu minimalnego $\alpha_i$
nad $K_{i-1}$ z przeniesionymi współczynnikami,
a takich różnych pierwiastków jest co najwyżej $d_i$.
Każdy wybór pierwiastka daje dokładnie jedno przedłużenie
na mocy Lematu~\ref{aa:lem:extend-embedding-simple}
(z izomorfizmem ciała pośredniego na jego obraz),
a różne pierwiastki dają różne przedłużenia.
Wynika stąd oszacowanie liczby zanurzeń przez
$\prod_i d_i=[L:K]$ na mocy twierdzenia o wieży.

Jeśli $L/K$ jest rozdzielcze, każdy $\alpha_i$ ma
rozdzielczy wielomian minimalny nad $K$. Wielomian
minimalny nad $K_{i-1}$ dzieli go w $K_{i-1}[X]$,
a więc także jest rozdzielczy. Po przeniesieniu
współczynników zachowuje stopień i brak pierwiastków
wielokrotnych. Ma więc dokładnie $d_i$ różnych
pierwiastków w $\Omega$ i na każdym etapie istnieje
dokładnie $d_i$ wyborów. Otrzymujemy $\prod_i d_i$
różnych zanurzeń.

Odwrotnie, załóżmy, że liczba wszystkich
$K$-zanurzeń $L\hookrightarrow\Omega$ wynosi $[L:K]$.
Dla dowolnego $\alpha\in L$ rozpocznijmy wieżę
od $K(\alpha)$, a następnie dołączmy skończenie
wielu pozostałych generatorów. Jeśli $m_{\alpha,K}$
ma mniej różnych pierwiastków niż jego stopień $d$,
na pierwszym etapie mamy mniej niż $d$ zanurzeń;
każde z nich przedłuża się do całego $L$ na co
najwyżej $[L:K(\alpha)]$ sposobów. Liczba
zanurzeń byłaby ściśle mniejsza niż
$d[L:K(\alpha)]=[L:K]$, co przeczy założeniu.
Zatem każdy $m_{\alpha,K}$ jest rozdzielczy.
\end{proof}

\begin{przyklad}[Dlaczego brak rozdzielczości zmniejsza liczbę zanurzeń]
\label{aa:prz:inseparable-embeddings}
Niech $K=\mathbb F_p(t)$ i $\alpha^p=t$. Rozszerzenie
$L=K(\alpha)$ ma stopień $p$, ponieważ $X^p-t$
jest nierozkładalny (Przykład~\ref{aa:prz:inseparable-example}).
Jednak $X^p-t=(X-\alpha)^p$ w $L[X]$, więc
każde $K$-zanurzenie $L\hookrightarrow\Omega$ musi
wysyłać $\alpha$ na $\alpha$. Jest zatem tylko jedno
takie zanurzenie, mimo że $[L:K]=p$.
\end{przyklad}

\section{Jednoznaczność ciała rozkładu}
\begin{twierdzenie}[Jednoznaczność ciała rozkładu z dokładnością do izomorfizmu]
\label{aa:tw:splitting-uniqueness}
Niech $\varphi:K\to K'$ będzie izomorfizmem ciał,
$f\in K[X]$ niezerowym wielomianem dodatniego stopnia,
$L/K$ ciałem rozkładu $f$ oraz $L'/K'$ ciałem rozkładu
$f^\varphi$. Wtedy istnieje izomorfizm $\Phi:L\to L'$
przedłużający $\varphi$. Nie musi on być jednoznaczny.
\end{twierdzenie}
\begin{proof}
Przeprowadzimy indukcję po $n=\deg f$. Dla $n=1$
oba ciała rozkładu są odpowiednio $K$ i $K'$, więc
wystarcza $\varphi$. Niech $n\ge2$ i przyjmijmy wynik
dla mniejszych stopni. Wybierzmy nierozkładalny czynnik
$g\in K[X]$ wielomianu $f$, a w $L$ pierwiastek
$\alpha$ wielomianu $g$. W $L'$ wybierzmy pierwiastek
$\beta$ wielomianu $g^\varphi$.
Lemat~\ref{aa:lem:extend-embedding-simple} daje
izomorfizm $\varphi_1:K(\alpha)\to K'(\beta)$
przedłużający $\varphi$ i wysyłający $\alpha$ na $\beta$.
W $K(\alpha)[X]$ mamy $f=(X-\alpha)h$,
gdzie $\deg h=n-1$. Po przeniesieniu współczynników
przez $\varphi_1$ otrzymujemy
$f^\varphi=(X-\beta)h^{\varphi_1}$.
Ciało $L$ jest ciałem rozkładu $h$ nad $K(\alpha)$:
rozszczepia $h$ i jest generowane nad $K(\alpha)$
przez pierwiastki $h$. Analogicznie $L'$ jest ciałem
rozkładu $h^{\varphi_1}$ nad $K'(\beta)$.
Założenie indukcyjne, stosowane do rozszerzenia bazowego
$K(\alpha)\cong K'(\beta)$, dostarcza izomorfizmu $L\to L'$
przedłużającego $\varphi_1$, a więc także $\varphi$.
\end{proof}

\section{Normalność i rozszerzenia Galois}
\begin{definicja}[Rozszerzenie normalne, rozdzielcze i Galois]
\label{aa:def:normal-separable-galois}
Niech $L/K$ będzie rozszerzeniem algebraicznym.
Nazywamy je \emph{normalnym}, jeśli każdy nierozkładalny
wielomian $f\in K[X]$ mający pierwiastek w $L$
rozkłada się w $L[X]$ na czynniki liniowe.
Rozszerzenie jest \emph{rozdzielcze}, jeśli dla każdego
$\alpha\in L$ wielomian minimalny $m_{\alpha,K}$
ma wyłącznie pierwiastki pojedyncze.
Rozszerzenie \emph{Galois} to rozszerzenie algebraiczne,
normalne i rozdzielcze; gdy dodatkowo $[L:K]<\infty$,
mówimy o \emph{skończonym rozszerzeniu Galois}.
\end{definicja}

\begin{twierdzenie}[Kryterium normalności przez zanurzenia]
\label{aa:tw:normal-embeddings}
Niech $L/K$ będzie rozszerzeniem skończonym i
$\Omega$ ciałem algebraicznie domkniętym zawierającym $L$.
Następujące warunki są równoważne:
\begin{enumerate}[label=\textup{(\roman*)}]
\item $L/K$ jest normalne;
\item dla każdego $K$-zanurzenia $\sigma:L\hookrightarrow\Omega$
      zachodzi $\sigma(L)=L$;
\item $L$ jest ciałem rozkładu pewnej skończonej rodziny
      niezerowych wielomianów z $K[X]$.
\end{enumerate}
\end{twierdzenie}
\begin{proof}
\emph{(i)$\Rightarrow$(ii).} Dla $\alpha\in L$ obraz
$\sigma(\alpha)$ zeruje $m_{\alpha,K}$.
Normalność gwarantuje, że wszystkie pierwiastki
tego wielomianu należą do $L$, stąd $\sigma(L)\subseteq L$.
Wtedy $\sigma:L\to L$ jest różnowartościowym
$K$-liniowym endomorfizmem przestrzeni wektorowej
skończonego wymiaru, więc jest surjektywny.

\emph{(ii)$\Rightarrow$(i).} Niech $f\in K[X]$ będzie
nierozkładalny i $\alpha\in L$ jego pierwiastkiem.
Dla dowolnego innego pierwiastka $\beta\in\Omega$
lemat o przedłużeniu zanurzeń daje $K$-zanurzenie
$K(\alpha)\to\Omega$ wysyłające $\alpha$ na $\beta$.
Przedłużamy je do $L$ na mocy
Wniosku~\ref{aa:wn:extend-embedding-finite}.
Z założenia jego obraz to $L$, zatem $\beta\in L$.
Wszystkie pierwiastki $f$ leżą w $L$, więc $f$
rozkłada się na czynniki liniowe w $L[X]$.

\emph{(i)$\Rightarrow$(iii).} Niech
$L=K(\alpha_1,\ldots,\alpha_r)$. Ponieważ $L/K$
jest normalne, każdy wielomian minimalny
$m_{\alpha_i,K}$ rozkłada się w $L[X]$.
Ciało generowane nad $K$ przez wszystkie pierwiastki
tych wielomianów zawiera wszystkie $\alpha_i$, a więc
jest równe $L$.

\emph{(iii)$\Rightarrow$(ii).} Każde $K$-zanurzenie
$\sigma:L\to\Omega$ wysyła pierwiastki odpowiednich
wielomianów na pierwiastki tych samych wielomianów.
Wszystkie one leżą w $L$, więc $\sigma(L)\subseteq L$;
argument wymiarowy z pierwszego kroku daje równość.
\end{proof}

\begin{twierdzenie}[Grupa Galois i stopień rozszerzenia]
\label{aa:tw:galois-group-degree}
Jeśli $L/K$ jest skończonym rozszerzeniem Galois, to
$\operatorname{Gal}(L/K):=\operatorname{Aut}_K(L)$
jest grupą skończoną oraz
$|\operatorname{Gal}(L/K)|=[L:K]$.
\end{twierdzenie}
\begin{proof}
Rozdzielczość i Twierdzenie~\ref{aa:tw:embeddings-count}
dają dokładnie $[L:K]$ różnych $K$-zanurzeń
$L\hookrightarrow\Omega$. Normalność oraz
Twierdzenie~\ref{aa:tw:normal-embeddings} pokazują,
że obraz każdego z nich jest równy $L$. Są to zatem
dokładnie wszystkie $K$-automorfizmy $L$.
\end{proof}

\begin{przyklad}[Rozszerzenie rozdzielcze, które nie jest normalne]
\label{aa:prz:separable-not-normal}
Dla rzeczywistego pierwiastka $\alpha=\sqrt[3]{2}$
rozszerzenie $\mathbb Q(\alpha)/\mathbb Q$ ma stopień trzy,
a więc jest skończone. W charakterystyce zero jest
rozdzielcze, ponieważ wszystkie nierozkładalne wielomiany
są rozdzielcze (Wniosek~\ref{aa:wn:charzero-separable}).
Nie jest normalne: nierozkładalny wielomian $X^3-2$
ma pierwiastek $\alpha$ w tym rozszerzeniu, lecz
pozostałe dwa pierwiastki są nierzeczywiste i nie należą
do $\mathbb Q(\alpha)\subseteq\mathbb R$.
Jedynym $\mathbb Q$-automorfizmem tego rozszerzenia
jest identyczność: obraz $\alpha$ musi być rzeczywistym
pierwiastkiem $X^3-2$, czyli $\alpha$.
\end{przyklad}

\begin{przyklad}[Rozszerzenie normalne, które nie jest rozdzielcze]
\label{aa:prz:normal-not-separable}
Niech $K=\mathbb F_p(t)$ i $L=K(\alpha)$, gdzie
$\alpha^p=t$. Najpierw uzasadnijmy $[L:K]=p$.
W pierścieniu $\mathbb F_p[t]$ wielomian $X^p-t$
traktujemy jako wielomian zmiennej $X$ o współczynnikach
w $\mathbb F_p[t]$. Jest on nierozkładalny nad
$\mathbb F_p(t)$ na mocy
Twierdzenia~\ref{aa:tw:eisenstein}
(kryterium Eisensteina), zastosowanego do dziedziny
euklidesowej $\mathbb F_p[t]$ i jej elementu pierwszego
$t$ (ideał $(t)$ jest pierwszy, gdyż iloraz to $\mathbb F_p$): współczynnik
wiodący jest równy $1$, wszystkie pozostałe
współczynniki nieprowadzące są podzielne przez $t$,
a wyraz wolny $-t$ nie jest podzielny przez $t^2$.
Dlatego $[L:K]=p$. Każdy element $x\in L$ ma $p$-tą potęgę w $K$.
Istotnie, dla $x=u(\alpha)/v(\alpha)$, gdzie
$u,v\in K[X]$ i $v(\alpha)\ne0$, endomorfizm
Frobeniusa daje
$x^p=u^{[p]}(t)/v^{[p]}(t)\in K$;
tutaj $u^{[p]}$ oznacza wielomian powstały przez
podniesienie współczynników $u$ do potęgi $p$.
Mianownik jest niezerowy, ponieważ jego wartość
wynosi $v(\alpha)^p$. Jeśli $m_{x,K}$
ma stopień większy od jednego, to dzieli
$X^p-x^p=(X-x)^p$ w $L[X]$. Wszystkie pierwiastki
$m_{x,K}$ są więc równe $x$, a zatem leżą w $L$:
rozszerzenie jest normalne. Nie jest rozdzielcze,
ponieważ wielomian minimalny $\alpha$ to $X^p-t$
i ma wielokrotny pierwiastek.
\end{przyklad}

\section{Ciała skończone jako rozszerzenia Galois}
\begin{twierdzenie}[Istnienie i jednoznaczność ciał skończonych]
\label{aa:tw:finite-fields-existence-uniqueness}
Dla każdej liczby pierwszej $p$ i każdej liczby całkowitej
$n\ge1$ istnieje ciało o $p^n$ elementach.
Każde dwa ciała o $p^n$ elementach są izomorficzne.
Ciało to oznaczamy $\mathbb F_{p^n}$. Jest ono ciałem
rozkładu $X^{p^n}-X$ nad $\mathbb F_p$.
\end{twierdzenie}
\begin{proof}
Niech $q=p^n$ i niech $L$ będzie ciałem rozkładu
$f=X^q-X\in\mathbb F_p[X]$, istniejącym na mocy
Twierdzenia~\ref{aa:tw:splitting-existence}. Pochodna
$f'=qX^{q-1}-1=-1$, więc $f$ ma dokładnie $q$
różnych pierwiastków w $L$. Oznaczmy ich zbiór przez $F$.
Dla $x,y\in F$ mamy $x^q=x$, $y^q=y$, a więc
$(x+y)^q=x^q+y^q=x+y$, ponieważ
podnoszenie do potęgi $p$ jest addytywne i iterujemy je
$n$ razy. Analogicznie $(xy)^q=xy$ oraz
$(-x)^q=-x$. Elementy $0,1$ należą do $F$,
a jeśli $x\in F\setminus\{0\}$, to
$(x^{-1})^q=(x^q)^{-1}=x^{-1}$.
Zbiór $F$ jest zatem podciałem $L$ i ma dokładnie
$q$ elementów. Zawiera $\mathbb F_p$, a jego elementy
są dokładnie wszystkimi pierwiastkami $f$, więc
$F$ jest ciałem rozkładu $f$.

Niech teraz $E$ będzie dowolnym ciałem o $q$ elementach.
Jego charakterystyka wynosi $p$: na mocy
Twierdzenia~\ref{aa:tw:finite-field-cardinality}
liczba elementów $E$ jest potęgą pewnej liczby
pierwszej $\ell=\operatorname{char}E$.
Równość $\ell^r=p^n$ i jednoznaczność rozkładu
na czynniki pierwsze wymuszają $\ell=p$.
Dla $x\in E^\times$ twierdzenie Lagrange'a daje
$x^{q-1}=1$, zatem $x^q=x$; to samo zachodzi dla $x=0$.
Wszystkie $q$ elementów $E$ są więc pierwiastkami $f$,
a $E$ jest ciałem rozkładu $f$ nad swoim kanonicznym
podciałem $\mathbb F_p$. Z jednoznaczności ciała
rozkładu (Twierdzenie~\ref{aa:tw:splitting-uniqueness})
otrzymujemy $E\cong F$.
\end{proof}

\begin{twierdzenie}[Automorfizmy ciała skończonego]
\label{aa:tw:finite-field-frobenius-group}
Niech $p$ będzie liczbą pierwszą, $n\ge1$ liczbą
całkowitą i $q=p^n$. Wtedy $\mathbb F_q/\mathbb F_p$
jest rozszerzeniem Galois stopnia $n$, a jego grupa
Galois jest cykliczna i generowana przez automorfizm
Frobeniusa $\operatorname{Fr}_p:x\mapsto x^p$.
\end{twierdzenie}
\begin{proof}
Ciało $\mathbb F_q$ jest ciałem rozkładu
$X^q-X$, a wielomian ten ma pochodną $-1$,
więc jest rozdzielczy. Rozszerzenie jest zatem
normalne i rozdzielcze; jego stopień wynosi $n$,
ponieważ jako przestrzeń wektorowa nad $\mathbb F_p$
ma $p^n$ elementów. Z
Twierdzenia~\ref{aa:tw:galois-group-degree} wynika,
że jego grupa Galois ma $n$ elementów.
Endomorfizm Frobeniusa jest automorfizmem
(Twierdzenie~\ref{aa:tw:field-frobenius}) i zachowuje
$\mathbb F_p$. Dla każdego $x\in\mathbb F_q$
mamy $\operatorname{Fr}_p^n(x)=x^q=x$, więc rząd
tego automorfizmu dzieli $n$. Jeśli $0<r<n$
i $\operatorname{Fr}_p^r$ był identycznością,
to wszystkie $q$ elementów byłyby pierwiastkami
niezerowego wielomianu $X^{p^r}-X$ stopnia
$p^r<q$, co jest niemożliwe. Rząd wynosi więc
$n$, czyli $\operatorname{Fr}_p$ generuje całą grupę.
\end{proof}

\begin{przyklad}[Ciało o ośmiu elementach]
\label{aa:prz:eight-field-frobenius}
Niech $F=\mathbb F_2[\alpha]$, gdzie
$\alpha^3+\alpha+1=0$. Wielomian $X^3+X+1$
nie ma pierwiastka w $\mathbb F_2$, więc jako
wielomian stopnia trzy jest nierozkładalny.
Mamy $[F:\mathbb F_2]=3$ i $|F|=8$.
Grupa Galois jest cykliczna rzędu trzy; jej
generatorami są $x\mapsto x^2$ oraz $x\mapsto x^4$.
Nie należy mylić tego faktu z cyklicznością
$F^\times$, która jest \emph{inną} grupą i ma rząd siedem.
\end{przyklad}

\begin{uwaga}[Droga do twierdzenia o odpowiedniości Galois]
\label{aa:uw:embeddings-next}
Wykazaliśmy, że w skończonym rozszerzeniu Galois
liczba automorfizmów jest równa stopniowi rozszerzenia.
W kolejnym rozdziale zbadamy, jakie podgrupy tych
automorfizmów odpowiadają ciałom pośrednim. Kluczowy
będzie dowód twierdzenia Artina o ciele elementów stałych,
bez zakładania z góry odpowiedniości Galois.
\end{uwaga}

\begin{uwaga}[Rola normalności w liczeniu automorfizmów]
Dla rozszerzenia $\mathbb Q(\sqrt[3]{2})/\mathbb Q$
stopień wynosi $3$, lecz jedynym automorfizmem nad $\mathbb Q$
jest identyczność: pozostałe dwa pierwiastki $T^3-2$ nie należą
do ciała rzeczywistego $\mathbb Q(\sqrt[3]{2})$.
Zatem równość liczby automorfizmów ze stopniem wymaga
założeń normalności i rozdzielczości, a nie samej skończoności.
\end{uwaga}
